% RESEARCH_DRAFT_SHA256: 9abf89d72cf6ae940e01c65fd44bfd4bb4cab79a4848d2dcf3145884bdbf419a
% Layout and source-reference wrappers expanded for the web renderer.
\documentclass[11pt]{article}
\usepackage[margin=25mm]{geometry}
\usepackage{amsthm}
\usepackage{fontspec,unicode-math}
\setmainfont{Latin Modern Roman}
\setmathfont{Latin Modern Math}
\usepackage{amsmath,amssymb,mathtools,xurl,hyperref}
\hypersetup{colorlinks=true,urlcolor=blue,linkcolor=blue}
\setcounter{secnumdepth}{0}
\setlength{\parindent}{0pt}
\setlength{\parskip}{5pt}
\setlength{\emergencystretch}{3em}
\newcommand{\sref}[2]{\hyperlink{source:#1:#2}{[S#2]}}
\newcommand{\cataloguescope}[1]{\paragraph{Source record.}#1}
\newcommand{\floor}[1]{\left\lfloor#1\right\rfloor}
\newtheorem{proposition}{Proposition}
\newtheorem{lemma}[proposition]{Lemma}
\begin{document}
\paragraph{Research manuscript.} OpenAI \texttt{gpt-6-astra}, reasoning \texttt{ultra}; 12 October 2026.
\paragraph{Preserved source body.} The complete catalogue mathematical text and sources follow with layout and source links expanded, including source qualifications and its historical preparation notes. New research begins at the section headed Research attempt.

\section{Evasive-LWE Research-Direction Discussion}\label{statement}

\subsection{Definitions and mathematical statement}

\paragraph{Shared notation.}Standard number systems are $\mathbb N,\mathbb Z,\mathbb Q,\mathbb R,\mathbb C$. Specialized notation and index ranges follow the cited source and the definitions below.


Learning with errors uses noisy linear equations in a finite algebraic domain. A private-coin algorithm keeps its random choices hidden; a public-coin setup makes designated randomness visible. Succinctness imposes a specified size bound.
\paragraph{Record scope and qualifications.}
The catalogue title is retained as metadata. Named certificates, contracts, bounds or frameworks in that title are not asserted as results of the original source extract.
The supplied registry flags this extraction as unrecoverable. It is displayed as a literal source fragment; the definitions provide context rather than a replacement formal statement.

The evasiveness notion, public/private randomness experiments, and desired succinctness parameters are absent; it is a research discussion.

Source-record limitation: Source flags the statement as unrecoverable: The official report records an open-ended discussion about which research direction to pursue after counterexamples, not a mathematically testable problem statement.
The source text below is rendered literally where necessary; damaged formulas and composite extracts are not silently repaired into a new mathematical claim.
The supplied catalogue marks this record as unrecoverable. Its extract and terminology are documented, but a complete original problem has not been reconstructed.
\paragraph{Statement recorded in the supplied source.}
Should we move away from private-coin evasive LWE to public-coin evasive LWE and succinct LWE? Discussion group: Which direction(s) should evasive LWE research take? Chris Brzuska \href{https://doi.org/10.4171/owr/2025/3}{[S2]}
\paragraph{Source status.}
The supplied catalogue labels this source record partially solved; that metadata was checked against the supplied snapshot on 2026-10-11. The label is preserved for provenance and is not an independent verification of present open status or mathematical correctness. The catalogue literature assessment is dated 2026-08-22.
\subsection{Short English statement}
The fragment discusses possible directions for evasive-learning-with-errors research.
\subsection{Sources}
\begin{itemize}
\item[S1] Ulam AI and the UnsolvedMath Contributors, UnsolvedMath: A Curated Collection of Open Mathematics Problems, catalogue entry. Catalogue curation and metadata: CC BY 4.0. The source extract is retained; definitions, English summaries and qualifications are editorial additions.
\url{https://huggingface.co/datasets/ulamai/UnsolvedMath}.
\item[S2] Cryptography (2025), Oberwolfach Reports, DOI 10.4171/owr/2025/3, pp. 141--174
\url{https://doi.org/10.4171/owr/2025/3}.
\item[S3] Brzuska, Uenal, and Woo, evasive-LWE attacks and variants
\url{https://eprint.iacr.org/2025/375.pdf}.
\item[S4] Lattice construction from public-coin evasive LWE
\url{https://eprint.iacr.org/2024/1416.pdf}.
\item[S5] OWR 2025/3 Cryptography report
\url{https://ems.press/content/serial-article-files/51349}.
\item Source-specific Oberwolfach Reports licences, including CC BY-SA 4.0 where applicable, are retained; see THIRD\_PARTY\_NOTICES.md.
\end{itemize}
\paragraph{Preparation and review.}
The original source and UnsolvedMath contributors retain attribution for the source extract. Codex (GPT-6; exact serving revision not exposed) assisted with definitions, English summaries, scope notes and formatting. No new human mathematical authorship or independently verified solution is asserted. Catalogue identity checked against the complete supplied registry and current repository. Distinct upstream records, including repeated questions and broad or damaged extracts, are retained for complete coverage of the source CS category. A repeated record is not a new independent conjecture. Complete mathematical-equivalence and current-literature review has not been performed. Preparation and final review follow contribution prompts C and D, with source gaps disclosed. Structural and rendering checks do not establish a new theorem or open status.
\label{end:statement}

\section{Research attempt}
\subsection{Literature, source context and scope}
The full original catalogue body and source qualifications are preserved above. The relevant primary workshop passages were read in the cached publisher report. Historical claims and subsequent attributed results are distinguished from the mathematical work below. A restricted example, standard lemma or known theorem is not counted as new progress on the entire original target. Missing assumptions and repeated or rhetorical extracts are not silently turned into different solved problems.

This preserved record is a discussion agenda after several attacks on different evasive-LWE formulations. The primary passage distinguishes private-coin samplers, samplers receiving a secret, public-coin variants and additional matrix inputs. Its historical account of known attacks is not a present-day blanket security assertion, and the extract supplies no precise new hardness experiment. A defensible limited attempt is to quantify one reason sampler information must be stated before comparing formulations.
\subsection{A finite leakage bound for adaptive evasion choices}
Let $R$ be any random variable, $\mathcal A$ an action set, and $E_a$ an event depending on $R$ for every fixed action $a$. Assume $\Pr[E_a]\le\varepsilon$ for all fixed $a$. Let deterministic leakage $L(R)$ take at most $2^\ell$ values, and let an adversary select $a=A(L(R))$. This is a finite information model, not an LWE assumption.
\begin{proposition}
The adaptive success probability is at most $\min(1,2^\ell\varepsilon)$. The factor $2^\ell$ can be attained exactly.
\end{proposition}
\begin{proof}
The range of the deterministic chooser $A$ contains at most $2^\ell$ actions. If it succeeds, at least one of the corresponding events $E_a$ occurs. A union bound gives $2^\ell\varepsilon$, and probability is at most one. If the adversary additionally uses independent private coins, condition on those coins and apply the same argument, then average.

For tightness take $R$ uniform on $\{0,1\}^n$, $0\le\ell\le n$, actions all $n$-bit strings, and $E_a=\{R=a\}$. Every fixed event has probability $\varepsilon=2^{-n}$. Reveal the first $\ell$ bits and choose them followed by $n-\ell$ zero bits. Success occurs precisely when the last $n-\ell$ bits are all zero, with probability $2^{\ell-n}=2^\ell\varepsilon$.
\end{proof}
\subsection{The exact failed inference}
A proof that every fixed test is rarely successful does not remain valid, with the same bound, when the test may depend on newly revealed sampler information. Conversely the union bound is not a cryptanalytic attack on an actual LWE variant: the events may be computationally hard to exploit, the leakage may have more structure, and the original hardness definition may already quantify over that information. The example gives neither an LWE instance nor a reduction between public- and private-coin formulations.
\section{Verification and remaining gaps}
Exhaustive strings for $n=6$ and $\ell=0,1,3,6$ verify the exact amplification. Selecting a surviving cryptographic assumption would require its full parameter distribution, sampler inputs, adversary order and success relation, followed by reductions or attacks. Those are absent from the agenda; this quantitative diagnostic does not manufacture them or recommend an unproved security claim.

\section{References and source attribution}
Primary workshop report: \url{https://doi.org/10.4171/owr/2025/3}. Relevant passages read: OWR 2025/3, printed pages155--157, Wu context and Brzuska discussion.

All retrievals were checked on 11 October 2026. Original question attribution remains in the preserved Sources section. External theorems are attributed and not claimed as new discoveries or fully verified merely by citation.


\paragraph{Integration review.} Codex (GPT-6; exact serving revision not exposed) reviewed the written arguments and expanded layout and source-reference wrappers. Repository metadata and the production web renderer were checked using Python 3.13 and JavaScript with MathJax 3.2.2. This model-assisted review is not independent human proof verification; no PDF compilation or formal proof checking is asserted.

\section{Completion Estimate}
\noindent\textbf{COMPLETION: 0\%}
No substantiated new progress toward the entire original target is claimed. The worked results and failed route are delimited above; a known result, illustrative model, or auxiliary obstruction is not counted as original-target completion.

\section{Final status and attribution}
\textbf{Attempt status: unresolved.} This describes this research attempt, not a new ruling on the source catalogue's status or on external literature claims.
The mathematics and this self-audit were generated by OpenAI \texttt{gpt-6-astra}, reasoning \texttt{ultra}, on 12 October 2026. Codex (GPT-6; exact serving revision not exposed) orchestrated the batch, read and reviewed the written mathematical arguments, independently reran the finite checks, and contributed formatting/repository review. No human mathematical authorship, independent human proof checking, or formal verification is claimed. The historical preparation paragraphs in the reproduced catalogue body concern that body only.

Contribution prompt A was applied to the complete definition and supplied source record, followed by the adversarial checks of prompt B. Local reproduction files are \texttt{batch026/checks.py}, \texttt{computations.json}, and the per-ID audit JSON. They were not public links at preparation time. Numerical checks are finite and do not prove unrestricted results. No LaTeX engine was found locally; compilation was not performed by this research agent.

\end{document}
